%% kapa: exam=1 round=22 year=2011 period=2 form=A subject=appraisal_law subject_ko=부동산관계법규 seq=22 number=62
%% source: Q-Net 감정평가사 기출문제 (cst00309) artlSeq=5036790 file=extracted/1008/5036790/2011년 제22회 감정평가사 1차시험 기출문제/2교시.hwp
%% conversion: kordoc 4.24.51 -> tools/gen1.py (auto); review=proofread-gpt-astra; flags=none
\begin{question}{62}
부동산 가격공시 및 감정평가에 관한 법령상 공동주택가격의 공시에 관한 설명으로 옳은 것은?
\begin{choices}
  \item 시장ㆍ군수 또는 구청장이 공동주택가격을 조사ㆍ산정하여 공시하여야 한다.
  \item 공동주택가격을 조사ㆍ산정하는 경우에는 둘 이상의 감정평가업자에게 이를 의뢰하여야 한다.
  \item 공동주택에 전세권이 설정되어 있는 경우 당해 전세권이 존재하지 아니하는 것으로 보고 적정가격을 조사하여야 한다.
  \item 공동주택가격은 매년 4월 30일을 공시기준일로 하여 산정ㆍ공시하여야 한다.
  \item 공동주택가격을 산정한 때에는 그 타당성에 대하여 감정평가업자의 검증을 받고 공동주택소유자등 이해관계인의 의견을 들어야 한다.
\end{choices}
\end{question}
\answer{3}
